% from ega0-1-fr.tex
\oldpage[0\textsubscript{I}]{12}
une suite exacte $A^p\to M\to 0$.
On dit qu'un $A$-module $M$ admet une \emph{présentation finie} s'il est isomorphe au conoyau d'un homomorphisme $A^p\to A^q$, autrement dit s'il existe une suite exacte $A^p\to A^q\to M\to 0$.
On notera que sur un anneau \emph{noethérien} $A$, tout $A$-module de type fini admet une présentation finie.

Rappelons qu'une $A$-algèbre $B$ est dite \emph{entière sur $A$} si tout élément de $B$ est racine dans $B$ d'un polynôme unitaire à coefficients dans $A$; il revient au même de dire que tout élément de $B$ est contenu dans une sous-algèbre de $B$ qui est un \emph{$A$-module de type fini}.
Lorsqu'il en est ainsi, et que $B$ est commutative, la sous-algèbre de $B$ engendrée par une partie finie de $B$ est un $A$-module de type fini; pour que l'algèbre commutative $B$ soit entière et de type fini sur $A$, il faut et il suffit donc que $B$ soit un $A$-module de type fini; on dit alors aussi que $B$ est une $A$-algèbre \emph{entière finie} (ou simplement \emph{finie} si aucune confusion n'en résulte).
On observera que dans ces définitions, on ne suppose pas que l'homomorphisme $A\to B$ définissant la structure de $A$-algèbre soit injectif.
\end{env}

\begin{env}[1.0.6]
\label{0.1.0.6-fr}
Un anneau \emph{intègre} est un anneau dans lequel le produit d'une famille finie d'éléments $\neq 0$ est $\neq 0$; il revient au même de dire que dans un tel anneau on a $0\neq 1$ et le produit de deux éléments $\neq 0$ est non nul.
Un idéal \emph{premier} d'un anneau $A$ est un idéal $\mathfrak{p}$ tel que $A/\mathfrak{p}$ soit intègre; cela entraîne donc $\mathfrak{p}\neq A$.
Pour qu'un anneau $A$ ait au moins un idéal premier, il faut et il suffit que $A\neq\{0\}$.
\end{env}

\begin{env}[1.0.7]
\label{0.1.0.7-fr}
Un anneau \emph{local} est un anneau $A$ dans lequel il existe un seul idéal maximal, qui est alors le complémentaire des éléments inversibles et contient tous les idéaux $\neq A$.
Si $A$ et $B$ sont deux anneaux locaux, $\mathfrak{m}$ et $\mathfrak{n}$ leurs idéaux maximaux respectifs, on dit qu'un homomorphisme $\varphi:A\to B$ est \emph{local} si $\varphi(\mathfrak{m})\subset\mathfrak{n}$ (ou, ce qui revient au même, si $\varphi^{-1}(\mathfrak{n})=\mathfrak{m}$).
Par passage aux quotients, un tel homomorphisme définit alors un monomorphisme du corps résiduel $A/\mathfrak{m}$ dans le corps résiduel $B/\mathfrak{n}$.
Le composé de deux homomorphismes locaux est un homomorphisme local.
\end{env}

\subsection{Racine d'un idéal. Nilradical et radical d'un anneau}
\label{subsection:0.1.1-fr}

\begin{env}[1.1.1]
\label{0.1.1.1-fr}
Soit $\mathfrak{a}$ un idéal d'un anneau $A$; la \emph{racine} de $\mathfrak{a}$, notée $\mathfrak{r}(\mathfrak{a})$, est l'ensemble des $x\in A$ tels que $x^n\in\mathfrak{a}$ pour un entier $n>0$ au moins; c'est un idéal contenant $\mathfrak{a}$.
On a $\mathfrak{r}(\mathfrak{r}(\mathfrak{a}))=\mathfrak{r}(\mathfrak{a})$; la relation $\mathfrak{a}\subset\mathfrak{b}$ entraîne $\mathfrak{r}(\mathfrak{a})\subset\mathfrak{r}(\mathfrak{b})$; la racine d'une intersection finie d'idéaux est l'intersection de leurs racines.
Si $\varphi$ est un homomorphisme d'un anneau $A'$ dans $A$, on a $\mathfrak{r}(\varphi^{-1}(\mathfrak{a}))=\varphi^{-1}(\mathfrak{r}(\mathfrak{a}))$ pour tout idéal $\mathfrak{a}\subset A$.
Pour qu'un idéal soit racine d'un idéal, il faut et il suffit qu'il soit intersection d'idéaux premiers.
La racine d'un idéal $\mathfrak{a}$ est l'intersection des idéaux premiers \emph{minimaux} parmi ceux qui contiennent $\mathfrak{a}$; si $A$ est noethérien, ces idéaux premiers minimaux sont en nombre fini.

La racine de l'idéal $(0)$ est encore appelée le \emph{nilradical} de $A$; c'est l'ensemble $\mathfrak{N}$ des éléments nilpotents de $A$.
On dit que l'anneau $A$ est \emph{réduit} si $\mathfrak{N}=(0)$; pour tout anneau $A$, le quotient $A/\mathfrak{N}$ de $A$ par son nilradical est un anneau réduit.
\end{env}

\begin{env}[1.1.2]
\label{0.1.1.2-fr}
Rappelons que le \emph{radical} $\mathfrak{R}(A)$ d'un anneau $A$ (non nécessairement commutatif) est l'intersection des idéaux à gauche maximaux de $A$ (et aussi l'intersection des idéaux à droite maximaux).
Le radical de $A/\mathfrak{R}(A)$ est $(0)$.
\end{env}

\subsection{Modules et anneaux de fractions}
\label{subsection:0.1.2-fr}

